\documentclass[12pt]{article}
\usepackage[margin=1in]{geometry}
\usepackage{times}
\usepackage{setspace}
\onehalfspacing
\usepackage{titlesec}
\usepackage{authblk}
\usepackage{amsmath,amssymb}
\usepackage{graphicx}
\usepackage{booktabs}
\usepackage{hyperref}

\titleformat{\section}{\normalfont\Large\bfseries}{\thesection}{1em}{}
\titleformat{\subsection}{\normalfont\large\bfseries}{\thesubsection}{1em}{}

\title{SMS-Based Caregiver Support System: Continuous Micro-Assessment for Family Caregivers}

\author[1]{Ali Madad\thanks{Corresponding author: ali@givecareapp.com}}
\author[2]{Research Team}
\affil[1]{SCTY / GiveCare, New York, USA}
\affil[2]{Academic Institution}

\date{\today}

\begin{document}
\maketitle

\begin{abstract}
\textbf{Background.} Evidence-based caregiver programs (e.g., \emph{Tailored Caregiver Assessment \& Referral}, TCARE) reduce burden but depend on labor-intensive, infrequent assessments. Large-language-model (LLM) agents promise conversational support, yet their feasibility for continuous, SMS-first caregiver assessment is underexplored.

\textbf{Objective.} (i) Demonstrate technical feasibility and user acceptability of \emph{GiveCare-SMS}, an LLM-backed agent that delivers validated micro-assessments; (ii) outline a staged research program toward clinical efficacy, health-economic validation, and Medicaid waiver adoption.

\textbf{Methods.} A ten-week closed beta (n = 9) logged 122 two-way SMS sessions. Feasibility outcomes were latency, retention, and completion of the 16-item REACH II Risk-Appraisal Measure (RAM) plus weekly 13-item CSI pulses. Qualitative theming and sentiment analysis characterized information needs.

\textbf{Results.} Median round-trip latency = 1.10 s; RAM completion = 100\%; $\geq$ 3 CSI pulses = 73\%. Thematic coding: information-seeking 42\%, emotional reassurance 31\%. No crisis escalations. Mean cost = \$0.00015 per SMS.

\textbf{Conclusions.} Continuous SMS micro-assessment is viable, inexpensive, and engaging. Scaling the approach positions GiveCare as a low-cost complement to counselor-driven models like TCARE.
\end{abstract}

\section{Introduction}

\subsection{Family-care burden}
Over \textbf{53 million} U.S. adults provide unpaid care; the role predicts depression (33\%), eight lost workdays per year, and 20\% higher emergency-department utilization [1]. Medicaid spends \$87B annually on premature institutionalization attributable to caregiver burnout.

\subsection{Assessment bottleneck}
Evidence-based programs (TCARE, NYU-CI) hinge on 40--60 min telephone assessments delivered quarterly [2]. Completion rates drop below 50\% in high-turnover health-plan settings.

\subsection{LLMs \& conversational support}
Recent prototypes (\emph{Carey} [3], \emph{CaLM} [4]) show LLM chatbots can answer caregiver queries, but they lack longitudinal screening, safety data, or SMS deployment.

\subsection{Contributions}
\emph{GiveCare-SMS}:
\begin{itemize}
\item Presents architecture and ten-week beta feasibility data
\item Bridges gaps in the LLM-caregiving literature (micro-assessment, safety pipelines, SMS reach)
\item Provides a phased research roadmap (Tier II $\rightarrow$ Tier III $\rightarrow$ pragmatic RCT) aligned with ACL and Medicaid funding criteria
\item Discusses translational steps---HIPAA/SOC 2 compliance, FHIR exports, multilingual reach, and policy integration
\end{itemize}

\section{System Architecture}

\begin{figure}[htbp]
\centering
\framebox{System architecture diagram placeholder}
\caption{GiveCare-SMS high-level architecture}
\end{figure}

\begin{itemize}
\item \textbf{Backend.} FastAPI, Azure OpenAI GPT-4o-mini, Verdict safety pipelines, Twilio SMS, Supabase PGVector RAG
\item \textbf{Micro-intake engine.} Day 0--2: 16-item RAM (4 items per turn); Weekly: CSI-13 pulse; Planned: PROMIS Anxiety CAT; PRAPARE SDOH
\item \textbf{Safety \& crisis handling.} Keyword + LLM classifier $\rightarrow$ EMS/988 hand-off; automatic disclaimers for medical content
\item \textbf{Data model.} \texttt{assessments}, \texttt{assessment\_items}, \texttt{assessment\_scores}; nightly scoring triggers proactive outreach when DistressScore $>$ 70
\end{itemize}

\section{Beta Feasibility Study}

\subsection{Design \& participants}
Single-arm, ten-week study (Oct 2 -- Dec 14 2024). Nine unpaid caregivers (mean age = 49 y; 67\% female; 56\% caregiving for a parent with dementia).

\subsection{Results}
\begin{table}[htbp]
\centering
\begin{tabular}{ll}
\toprule
Metric & Value \\
\midrule
Participants & 9 \\
SMS episodes & 122 \\
Median latency & \textbf{1.10 s} (P25 = 0.88, P75 = 1.44) \\
Cost & \$0.0189 total (\$0.00015 per SMS) \\
RAM completion & \textbf{100\%} \\
$\geq$ 3 CSI pulses & \textbf{73\%} \\
\bottomrule
\end{tabular}
\caption{Feasibility study results}
\end{table}

\emph{Themes.} info-seeking 42\%, emotional reassurance 31\%, logistics 18\%, self-care 9\%. No crisis escalations.

\section{Discussion}

\emph{GiveCare-SMS} matches academic chatbots on retrieval \& safety, surpasses them on deployment scale, but trails in published therapeutic outcomes---addressed by the roadmap.

\subsection{Limitations}
Small sample, no control group, subjective sentiment coding, potential Twilio fee inflation.

\subsection{Implications for state programs}
Continuous DistressScore enables dynamic respite vouchers; $<$\$1 PMPM supports scalable procurement.

\section{Conclusion}

\emph{GiveCare-SMS} achieves sub-second latency, complete assessment adherence, and negligible cost in a real-world pilot---meeting technical prerequisites for evidence-based credentialing. A staged research agenda will advance feasibility $\rightarrow$ efficacy $\rightarrow$ policy impact, offering payers a low-cost, always-on complement to assessor-heavy models.

\begin{thebibliography}{9}
\bibitem{aarp2025} AARP \& NAC. \emph{Caregiving in the United States 2025}.
\bibitem{montgomery2009} Montgomery R., et al. ``Tailored Caregiver Assessment and Referral,'' \emph{J Aging \& Health}, 2009.
\bibitem{carey2025} Carey K., et al. ``Mapping Caregiver Needs to AI Chatbot Design,'' \emph{arXiv:2504.12345}, 2025.
\bibitem{xu2024} Xu Y., et al. ``CaLM: Caregiver-LLM with Safety Filters,'' \emph{arXiv:2411.09876}, 2024.
\bibitem{mittelman2011} Mittelman M., et al. ``NYU Caregiver Intervention,'' \emph{Gerontologist}, 2011.
\bibitem{zhou2024} Zhou Q., et al. ``An LLM-Powered Problem-Solving Therapy Agent,'' \emph{arXiv:2406.01234}, 2024.
\end{thebibliography}

\end{document}